\documentclass[11pt,a4paper]{article}
\usepackage[margin=2.5cm]{geometry}
\usepackage{tabularray}
\UseTblrLibrary{booktabs}
\DefTblrTemplate{firsthead, middlehead, lasthead}{default}{}
\usepackage{amsmath,amssymb,amsthm}
\usepackage{xcolor}
\usepackage{hyperref}
\usepackage{enumitem}
\usepackage{booktabs}
\usepackage{array}
\usepackage{microtype}
\usepackage{parskip}

\definecolor{sectionblue}{RGB}{0,70,140}
\definecolor{calloutgray}{RGB}{240,240,245}
\definecolor{warningamber}{RGB}{200,120,0}
\definecolor{proofgreen}{RGB}{0,110,60}

\hypersetup{
  colorlinks=true,
  linkcolor=sectionblue,
  urlcolor=sectionblue,
  citecolor=sectionblue
}

\newcommand{\vp}{\varphi}
\newcommand{\HH}{\mathbb{H}}
\newcommand{\CC}{\mathbb{C}}
\newcommand{\OO}{\mathbb{O}}
\newcommand{\ZZ}{\mathbb{Z}}
\newcommand{\aem}{\alpha_{\mathrm{em}}}
\newcommand{\twoI}{2I}
\newcommand{\Ehat}{\hat{E}_8}
\newcommand{\bst}{\beta^*}
\newcommand{\ERy}{E_{\mathrm{Ry}}}
\newcommand{\TCMB}{T_{\mathrm{CMB}}}
\newcommand{\Tsirelson}{2\sqrt{2}}
\newcommand{\CHSH}{\mathrm{CHSH}}
\newcommand{\xcond}{x_{\mathrm{cond}}}
\newcommand{\vEW}{v_{\mathrm{EW}}}
\newcommand{\mH}{m_H}

\newenvironment{callout}[1]{%
  \begin{trivlist}\item
  \colorbox{calloutgray}{\parbox{0.96\textwidth}{%
  \smallskip\textbf{#1}\par\smallskip
}{}}%
}{\end{trivlist}}

\newenvironment{prerequisite}{%
  \begin{trivlist}\item
  \colorbox{calloutgray}{\parbox{0.96\textwidth}{%
  \smallskip\textit{QFT/GR prerequisites for this section:}\par\smallskip
}{}}%
}{\end{trivlist}}

\title{\textbf{A Reader's Guide to the Density-Feedback Hopfion Series\\[4pt]
\large Papers I--XVII: Entry Points, Technical Prerequisites, and Critical Evaluation}}
\author{A technical companion for readers with background in QFT and GR}
\date{2026}

\begin{document}
\maketitle
\tableofcontents
\newpage

%═══════════════════════════════════════════════════════════════════════════════
\section{Purpose and scope of this guide}
%═══════════════════════════════════════════════════════════════════════════════

This guide is written for a reader with graduate-level training in quantum field
theory and general relativity who wants to evaluate the density-feedback
Faddeev--Niemi Hopfion series (Papers~I--XVII) critically. It does three things.

First, it maps the logical structure of the series: which results follow from
which, what is proved as a theorem, what is asserted as a claim, and where the
genuine open questions are. Second, it identifies the minimum QFT/GR background
needed to enter each paper without losing the thread. Third, it explicitly
addresses the question a QFT-trained reader will ask first: does this framework
contradict anything that is experimentally established, and if not, why not?

The short answer to that last question is: no, it does not. The framework
re-derives, from a single-input condensate model, the same numerical values that
the Standard Model accommodates as free parameters. It does not predict
deviations from known SM observables at any energy scale that has been probed;
it predicts deviations only where the SM is silent (the cosmological sector,
lepton Bell tests, and BAO peak position). The reader who expects a contradiction
with established particle physics will not find one, and this guide explains
precisely why.

A word on the Foundational Companion paper, which is an ideal first read for
anyone approaching the series. That paper collects three independent proofs that
the WZW level is $k=3$, shows that every appearance of $\vp$ in the series is a
consequence of the single equation $k+2=5$, and establishes the three-generation
lepton count as a hard algebraic theorem. It is self-contained in about nine pages
and requires only familiarity with WZW models and the Verlinde formula. Read it
first.

%═══════════════════════════════════════════════════════════════════════════════
\section{The condensate in brief: what kind of theory is this?}
%═══════════════════════════════════════════════════════════════════════════════

\subsection{The field and its action}

The theory has a single real scalar field $\rho(x)$, the condensate density, governed
by a $k$-essence action:
\begin{equation}
  S_{\mathrm{cond}} = \int d^4x\,\sqrt{-g}\,\frac{|\nabla\rho|^2}{\vp^6(1+\bst|\nabla\rho|^2)},
  \qquad \bst = 0.452.
  \label{RG:eq:action}
\end{equation}
This is not an ad hoc choice. Paper~I~\cite{Paper1} establishes that the saddle-point solution of
the density-feedback Faddeev--Niemi functional (a baby-Skyrme type functional on
maps $S^3\to S^2$ with anisotropic gradient coupling) forces $\vp^6$ in the
denominator through the Bogomolny self-duality condition $Q_{\mathrm{num}}=Q_{\mathrm{group}}$,
where $Q=10$ is the topological Hopf charge. The coupling $\bst=0.452$ is fixed
by a second self-consistency condition (the feedback stiffness $K_{\mathrm{fb}} = 2\vp J_a$),
not by fitting to data. Paper~II~\cite{Paper2} proves the saddle-point condition analytically via a Pohozaev--Derrick identity.

A QFT reader will recognise~\eqref{RG:eq:action} as a DBI-type or Born--Infeld kinetic
term, closely related to $k$-essence dark energy models. The key difference from
standard $k$-essence is that $\vp^6$ and $\bst$ are not free parameters: they are
uniquely determined by the topology of the Hopfion solution.

\subsection{What the condensate is not}

The condensate is not a new fundamental particle added to the Standard Model.
It is identified as the ground-state configuration of the existing scalar sector:
the gradient-energy condensate of a topological soliton (a Hopfion with charge
$Q=10$ in the icosahedral representation of $\mathrm{SU}(2)_3$) that is present
in the vacuum of a scalar-tensor theory naturally consistent with gravity.
Standard Model particles arise as
excitations of this condensate, not as separate fields coupled to it.

The QFT reader should think of this as analogous to how skyrmions describe
baryons in the Skyrme model: the topological structure of the field configuration,
not any particular interaction term, produces the observed quantum numbers.

\subsection{One input and it's role}

The entire framework takes one measured inputs:
\begin{equation}
  \TCMB = 2.7255\,\text{K}.
\end{equation}
$\TCMB$ sets the condensate scale $\mathcal{L}_{\mathrm{cond}}$ via
$\mathcal{L}_{\mathrm{cond}} = \hbar c / (k_B \TCMB)$ (the de Broglie
wavelength at the CMB temperature).
Everything else is derived.
Paper~XII~\cite{Paper12} (the profile normalisation paper) closes the last logical gap:
it proves that the condensate's gradient energy density equals the CMB
photon energy density at thermodynamic equilibrium, which forces
$m_e/\mathcal{L}_{\mathrm{cond}} = \vp^{20}\,e^{4\pi-3/400 -\alpha/\pi}$
(residual $0.013\%$).
The framework therefore reduces to \emph{one} measured input, $\TCMB$,
with $m_e$ becoming a derived quantity.

This extreme parsimony is the framework's central claim and its central
vulnerability: if any single derived prediction is wrong, the framework fails
as stated. Section~\ref{RG:sec:consistency} catalogues the predictions and their
current observational status.

%═══════════════════════════════════════════════════════════════════════════════
\section{Reading map: how the papers relate}
%═══════════════════════════════════════════════════════════════════════════════

\begin{center}
\begin{longtblr}[
  caption = {},
  label = {tab:paper_relations},
] {
  colspec = {X[1,p] X[3.5,p] X[8.5,c]}, % Proportional widths
  rowhead = 1, % Repeat header on each page
  %row{1} = {font=\bfseries},
}
\toprule
Paper & Core result & Depends on \\
\midrule
I & Saddle-point, $Q=10$, $k=3$, Koide relation &
  Faddeev--Niemi functional; $\twoI$ McKay correspondence \\
II & Analytic proof of LSC via Pohozaev identity &
  Variational calculus; Derrick scaling \\
III & LSC closed; $\aem$ two-term formula ($0.004\%$) &
  Witten non-abelian bosonization; $\mathrm{SU}(2)_3$ WZW \\
IV & Seven extensions: $\aem$ three-term ($0.0004\%$) and
  four-term ($5\times10^{-7}\%$); spoke spectrum;
  mass renormalization &
  One-loop WZW; Chern--Simons; Verlinde formula \\
V & Strong-force sector: QCD coupling, quark masses,
  CKM matrix, CP violation &
  $\mathbb{CP}^2$ Hopf fibration; semi-Dirac dispersion;
  $E_8$ Coxeter exponents \\
VI & Golden-ratio mass tower from RG fixed point;
  lepton mass formula &
  Derrick scaling; $k$-essence RG; WZW T-matrix \\
VII & Gravity from condensate; GR; inflation; dark energy; dark matter;
  cosmological constant &
  Scalar-tensor EFT; chameleon mechanism; Starobinsky inflation \\
VIII & Bell inequality geometry; CHSH from $\sin^4\theta/\vp^6$;
  Paper~IX~\cite{Paper9} Monte Carlo &
  Quantum information; Bell inequalities \\
IX & Born rule from $\HH\to\CC$ projection; Tsirelson bound;
  Malus interaction Hamiltonian &
  Division algebras; Hopf fibration; Fourier analysis on $S^1$ \\
X & Canonical quantisation; Hilbert space; golden-ratio spectrum;
  Bekenstein--Hawking &
  Geometric quantisation; Dirac constraint; spectral zeta function \\
XI & $\vp$-spiral periodic table; Rydberg as tower level;
  $Z_{\mathrm{eff}}(\mathrm{O})=2$ proved &
  Corollary of Papers~III and~IV; HF atomic theory \\
XII & Profile normalisation conjecture proved; $m_e$ derivable from $\TCMB$ alone &
  Papers~I--III (energy integrals, virial); thermodynamic equilibrium axiom \\
XIII & Atomic QM from condensate; Schr\"odinger equation, Pauli exclusion,
  spin-$\tfrac{1}{2}$, mass anisotropy &
  Paper~X~\cite{Paper10} (quantisation); Paper~VII~\cite{Paper7} (chameleon); Paper~I~\cite{Paper1} (Hopf charge) \\
XIV & Thick-torus toroidal correction; Yukawa closure;
  $\delta\vEW/\vEW$ from $0.69\%$ to $0.04\%$ &
  Papers~I--III (profile integrals); Paper~IV~\cite{Paper4} (Yukawa formula) \\
XV & $Q_H=3$ trefoil sector; quark charges from conformal weights;
  $\lambda_3=\vp^6$ proved; normalisation identity $\vp^9\sqrt{J_{2a}J_4}=3$ &
  Papers~II--III (LSC, bosonization); McKay correspondence;
  $\mathrm{SU}(3)_1$ WZW; Verlinde formula; fusion categories \\
XVI & Quark generation mass hierarchy: systematic survey of six ruled-out
  mechanisms; Y-junction string-tension energy isolated and computed;
  one-loop multi-threshold QCD running of the colour-shift prediction &
  Paper~XV (trefoil sector, colour shift $\delta n^{(C)}=\tfrac12$);
  Paper~VI (tower incommensurability); Paper~V ($E_8$ Coxeter exponents) \\
XVII & $Q_H=1$ sector identified as the neutrino: unknotted topology,
  $\Ztwo\leftrightarrow\widehat{A}_1$ McKay correspondence, $\mathrm{SU}(2)_1$
  WZW, colour exclusion strengthened; parameter-free mass prediction
  $m_{\nu_1}\approx42\,\mathrm{meV}$; PMNS matrix as an exact permutation &
  Papers~XV--XVI (sector-identification method); Paper~VI (tower
  irrationality); Paper~I (McKay correspondence, Hopf charge) \\
\bottomrule
\end{longtblr}
\end{center}

The logical spine is I$\to$III$\to$VI$\to$IX$\to$X$\to$XIII$\to$XV. Papers II and IV are
extensions of I and III respectively; Papers V and VII branch off to the
strong and gravitational sectors; Paper VIII~\cite{Paper8} provides numerical verification for IX;
Paper XII closes the proof of the one-parameter chain;
Paper~XIII~\cite{Paper13} extends the quantum theory to atomic physics;
Paper~XIV~\cite{Paper14} closes the electroweak residual;
Paper~XV~\cite{Paper15} extends the framework to the $Q_H=3$ quark sector
and proves the trefoil Bogomolny parameter $\lambda_3=\vp^6$;
Paper~XVI~\cite{Paper16} surveys the $Q_H=3$ generation-mass problem opened
by Paper~XV, ruling out six candidate mechanisms and isolating the
Y-junction string-tension energy as the remaining candidate source of
generation-dependence;
Paper~XVII~\cite{Paper17} completes the sector inventory by identifying
$Q_H=1$ as the neutrino sector and derives a parameter-free neutrino mass
prediction testable by Project~8.
A reader who wants the cleanest path to the most novel result should read the
Foundational Companion, then I, III, VI, IX, X, XII, XIII, XIV, XV, XVII in order;
XVI is best read as a companion to XV for readers specifically interested in
the quark-mass open problem.

%═══════════════════════════════════════════════════════════════════════════════
\section{Paper-by-paper entry guide}
%═══════════════════════════════════════════════════════════════════════════════

\subsection{Foundational Companion (start here)}

\textit{Prerequisites:} $\mathrm{SU}(2)_k$ WZW models; Verlinde formula;
McKay correspondence; binary icosahedral group $\twoI$.

The paper proves one main theorem with five items. Items (a)--(c) give three
independent derivations of $k=3$:
\begin{itemize}[noitemsep]
\item \textbf{Route 1 (numerical):} The formula
  $\aem^{-1} = 360/\vp^2 - k/(2\pi) + 1/(9\vp^6) - 1/(36\pi\vp^6)$ from Paper~IV~\cite{Paper4}
  gives $137.036$ at $k=3$ and excludes $k=4$ at $1.5\times10^7\,\sigma$.
\item \textbf{Route 2 (algebraic):} The quantum dimension $d_{1/2} = 2\cos(\pi/(k+2))$
  equals $\vp = 2\cos(\pi/5)$ if and only if $k+2=5$, i.e.\ $k=3$.
  This is a bijective proof.
\item \textbf{Route 3 (algebraic, strongest):} The WZW truncation
  $j \leq k/2$ allows exactly three non-vacuum primaries at $k=3$
  and algebraically excludes $j=2$, hence a fourth lepton generation.
  This is independent of any numerical input.
\end{itemize}
Items (d)--(e) add the quark sector: lepton generations use SU(2)$_3$
primary truncation; quark generations use $E_8$ Coxeter exponent saturation
(six of eight exponents occupied; $\{23,29\}$ predict no fourth quark generation).
These are structurally distinct but both rooted in the same $\twoI/\Ehat$ structure.

The reader who finishes this paper knows the algebraic skeleton of the series.

\subsection{Paper I: The Hopfion saddle point}

\textit{Prerequisites:} Faddeev--Niemi functional; Hopf fibration $S^3\to S^2$;
homotopy theory $\pi_3(S^2)=\ZZ$; McKay correspondence for finite subgroups of
$\mathrm{SU}(2)$.

Paper~I~\cite{Paper1} establishes that the saddle-point solution of the density-feedback
functional has topological charge $Q=10$, proves the McKay correspondence
$\twoI\leftrightarrow\Ehat$, and derives the WZW level $k=3$ from
$\mathrm{ord}(q_{\twoI}) = 2(k+2) = 10$. This is the foundational algebraic
identification from which all subsequent results flow.

Key subtlety for the QFT reader: the Hopfion charge $Q=10$ is not an input.
It emerges from the discrete icosahedral symmetry of the saddle-point orbit
under $\twoI\subset\mathrm{SU}(2)$, and from the McKay correspondence.
The Koide relation for lepton masses is derived as a corollary of the same
$\twoI$ representation theory.

\subsection{Papers II--III: The analytic backbone}

\textit{Prerequisites (II):} Pohozaev--Derrick identity; variational analysis
of nonlinear elliptic equations.

\textit{Prerequisites (III):} Witten non-abelian bosonization ($\mathrm{SU}(N)_k$
WZW from $N$ free Dirac fermions); Verlinde formula for fusion multiplicities;
modular S-matrix.

Paper~III~\cite{Paper3}contains the most important single result for a skeptical reader:
the identification of the density-feedback condensate as the $\mathrm{SU}(2)_3$
WZW model of three free Dirac fermions (Witten bosonization, Theorem
~10.1). It is an exact equivalence
proved via the LSC (Linking Scale Conjecture, proved as Theorem~2.2).
The $\alpha$ two-term formula follows from the Verlinde fusion rules as a direct
consequence: the golden angle $360/\vp^2$ is the geometric attractor of the
icosahedral condensate, and $k/(2\pi)$ is the WZW anti-screening correction.

A QFT reader will note that Witten bosonization normally applies in $1+1$ dimensions.
Paper~III~\cite{Paper3} proves the full 3+1D equivalence in two steps.
First, Theorem~10.3 establishes via
explicit Kaluza--Klein reduction that the 3+1D $k$-essence action reduces to the
2D WZW zero mode at energies $E \ll r_{\mathrm{core}}^{-1}$, with corrections
bounded by $O(R_0^{-2})$; the Wess--Zumino term descends correctly via the
Zumino--Faddeev descent equations; and the Hopf charge $Q=10$ is preserved by
a homotopy argument.
Second, Theorem~10.1 identifies the 2D zero-mode theory
as the $\mathrm{SU}(2)_3$ WZW model via the Pohozaev identity and Witten
non-abelian bosonization.
The $\alpha$ formula then follows from the $\mathrm{SU}(2)_3$ Verlinde fusion rules.

\subsection{Paper IV: The three and four-term $\aem^{-1}$ formula and the mass spectrum}

\textit{Prerequisites:} One-loop WZW mass renormalisation; Chern--Simons amplitude;
Verlinde formula; modular T-matrix.

Paper~IV~\cite{Paper4} extends Paper~III~\cite{Paper3} in seven directions, of which the most significant
for an independent reader are: (i) the three-term $\alpha$ formula
($\aem^{-1} = 360/\vp^2 - k/(2\pi) + 1/(9\vp^6) = 137.036\,49$, $0.0004\%$ accuracy,
Theorem~5.1), which is a one-loop WZW correction to the
leading-order result of Paper~III; the four-term extension
($\aem^{-1} = 360/\vp^2 - k/(2\pi) + 1/(9\vp^6) - 1/(36\pi\vp^6) = 137.036\,00$,
$5\times10^{-7}\%$ accuracy, Theorem~5.3),
where the fourth term is the Chern--Simons correction to the QED vacuum
polarisation inside the condensate background; and (ii) the spoke mass spectrum
(Theorem~6.11), which gives the ratio $m_e/\mathcal{L}_{\mathrm{cond}}
= \vp^{20}\cdot e^{4\pi}\cdot(1+O(R_0^{-2}))$ from the WZW Chern--Simons
one-loop amplitude.

The reader should verify both formulas directly.
The three-term formula uses $\vp^2 = (3+\sqrt{5})/2$, so $360/\vp^2 = 720/(3+\sqrt{5})$;
$k/(2\pi) = 3/(2\pi)$; and $1/(9\vp^6) = 1/(9((1+\sqrt{5})/2)^6)$.
The fourth term adds $-1/(36\pi\vp^6)$: exactly $T_3/(4\pi)$ with a sign flip,
where the $4\pi$ is the Chern--Simons coupling normalisation at unit WZW level.
All four computations are elementary and reproducible.

Two points about the four-term series are worth holding clearly.
First, the two-loop WZW correction is \emph{exactly zero}: the WZW beta
function is one-loop exact (the model is a CFT), and the higher primaries
$j=1$ and $j=3/2$ both satisfy $\cos(2\pi Q T_j) = 0$ at $k=3$, $Q=10$,
contributing nothing to the real part of the T-matrix sum.
The series is genuinely complete at four terms within the WZW framework.
Second, the convergence ratio $T_4/T_3=1/(4\pi)\approx0.08$; Paper~IV
estimates a further term at this rate would be $\sim6\times10^{-6}$ in
$\aInv$ (Paper~IV, Remark 5.4), for illustration of convergence only ---
the paper states the WZW series itself terminates at four terms
(two-loop WZW vanishes exactly, as above). The current four-term
residual, $|137.035\,998\,486-137.035\,999\,177|=6.9\times10^{-7}$, is
well within experimental reach (current precision
$\sim1$--$3\times10^{-8}$, e.g.\ Morel et al.\ 2020); Paper~IV
attributes it to ordinary sub-leading QED physics outside the
condensate framework with an as-yet-unverified candidate closed form.

\subsection{Paper V: The strong force sector}

\textit{Prerequisites:} $\mathbb{CP}^2$ Hopf fibration $S^1\hookrightarrow S^5\to\mathbb{CP}^2$;
semi-Dirac anisotropic dispersion; $E_8$ root system and Coxeter exponents;
CKM matrix; renormalisation group running of $\alpha_s$.

Paper~V~\cite{Paper5} is the most technically involved paper in the series. It derives the
QCD coupling $\alpha_s(M_Z)$, the quark mass ratios, the CKM matrix, and the CP
violation phase from the same condensate structure, using the $\mathrm{SU}(3)$
colour gauge group arising from the homotopy group $\pi_5(\mathbb{CP}^2)=\ZZ$.

The central claim that requires independent scrutiny: the effective spatial
dimension $d_{\mathrm{eff}}^{\mathrm{exact}} = 43/14$ (Proposition~5.1).
This arises from the semi-Dirac anisotropic dispersion ($k_x^2/(2m_*)$ in one
direction, $|k_y|$ in the perpendicular direction), which gives a density of
states exponent $d_{\mathrm{eff}} = 1/2 + 3/2 = 2$ at leading order corrected
to $43/14$ at subleading order. The reader should check this via the explicit
dispersion integral (equation 9 of Paper~V~\cite{Paper5}).

\subsection{Paper VI: The mass tower and electroweak scale}

\textit{Prerequisites:} Derrick scaling; RG fixed-point analysis;
$k$-essence Lagrangian; WZW T-matrix phases.

Paper~VI~\cite{Paper6} derives the fermion mass formula from first principles. The golden ratio
is not assumed; it emerges as the \emph{unique} RG fixed point of the condensate
$k$-essence action under Derrick scaling (Theorem~):
\begin{equation}
  \frac{d\mathcal{E}[\rho_\lambda]}{d\lambda}\bigg|_{\lambda=\vp} = 0,
  \quad
  \frac{d^2\mathcal{E}[\rho_\lambda]}{d\lambda^2}\bigg|_{\lambda=\vp} > 0,
\end{equation}
where $\rho_\lambda(x) = \lambda^d\rho(\lambda x)$ is the Derrick dilation.
The mass tower $m_n = m_0\,\vp^{-2n}$ follows as the spectrum of stable bound
states at the fixed point. The WZW T-matrix phases then place each lepton at a
precise irrational position between tower rungs via the two-clock structure
(Sections 4--5 of Paper~VI).

\subsection{Paper VII: Gravity, cosmology, inflation}

\textit{Prerequisites:} Scalar-tensor gravity; Brans--Dicke theory; chameleon
mechanism; Vainshtein screening;
Starobinsky inflation; $k$-essence dark energy; Cassini PPN constraint.

Paper~VII~\cite{Paper7} asks: what does the condensate action~\eqref{RG:eq:action}, now in $3+1$
dimensions with non-minimal curvature coupling, predict for cosmology?
The effective scalar-tensor action in the Jordan frame is
\begin{equation}
  S_{\mathrm{eff}} = \int d^4x\,\sqrt{-g}\,
  \Bigl[\frac{F(\rho)}{2}R - \frac{|\nabla\rho|^2}{\vp^6(1+\bst|\nabla\rho|^2)}\Bigr],
  \label{RG:eq:Seff}
\end{equation}
where $F(\rho_\infty) = M_{\mathrm{Pl}}^2(39\vp+25)/2$ at the attractor density
(an exact golden-ratio identity from $\vp^2=\vp+1$, Theorem \texttt{thm:einstein}).
Energy-momentum conservation $\nabla^\mu T_{\mu\nu}=0$ is derived, not assumed
(Theorem \texttt{thm:conservation}).

\textbf{Scalar-tensor and chameleon structure.}
The denominator $1/(1+\bst\rho)$ in the kinetic term plays a dual role.
At \emph{cosmological} densities $\rho\sim\rho_\infty=\vp/\bst$:
the effective coupling $\beta_{\mathrm{eff}}(\rho) = \bst/(1+\bst\rho)^2$
is suppressed, screening the scalar from gravitational coupling at large scales
(the chameleon mechanism (Khoury--Weltman 2004)).
At \emph{Solar System} densities $\rho\gg\rho_\infty$:
the nonlinear kinetic term produces a large self-coupling that suppresses
scalar fifth forces inside the Vainshtein radius
$r_V \approx 9.3\times10^6\,\mathrm{AU}$ (Theorem~6.9,
Corollary~6.12).
The Vainshtein radius is a \emph{topological} invariant:
$r_V \propto (G_N M_\odot / m_\xi^2)^{1/3}$
with $m_\xi \sim H_0$ set by the condensate mass, which in turn is fixed
by $Q=10$ and $\rho_\infty = \vp/\bst$.
These two mechanisms operate at different scales and are not competing:
the chameleon accounts for screening at $r\gg r_V$; the Vainshtein
radius suppresses fifth forces at $r\ll r_V$.
Remark~\texttt{rem:chameleon} discusses this complementarity explicitly;
a QFT reader should work through it before accepting Paper~VII's
Cassini resolution.

\textbf{The cosmological constant.}
The vacuum energy is not absent but bound inside the condensate. The cosmological constant problem is addressed via the chameleon mechanism
built into the density-feedback denominator $1/(1+\bst\rho)$: the vacuum energy
is screened at cosmological densities to give  $\Lambda_{\mathrm{obs}} \approx 10^{-122}M_{\mathrm{Pl}}^4$ as a geometric consequence of $\rho_\infty = \vp/\bst$ and $1+\vp=\vp^2$. The relevant attractor density $\rho_\infty$ is a fixed point of the
condensate equation of motion; the screening is not a fine-tuning but
a saddle-point condition.

\textbf{Dark matter.}
Condensate perturbations cluster gravitationally with sound speed
$c_s = 1/\vp$ at density $\rho_\infty$ (exact from $\vp^2=\vp+1$).
The Jeans length is $\approx20{,}000\,\mathrm{Mpc}$, so the condensate
is pressureless on all galactic and cluster scales. No new particle is introduced.
The primary prediction distinguishing this from CDM:
a BAO peak shift by factor $\vp$ to $k_{\mathrm{BAO}}^{\mathrm{cond}}
\approx0.0346\,\mathrm{Mpc}^{-1}$ (falsifiable with DESI DR2, Euclid).

\textbf{Inflation.}
The non-minimal coupling of $\rho$ to curvature gives a Higgs--Starobinsky
branch with $n_s=0.9654$ ($0.12\sigma$ from Planck 2018),
$r=0.0036$ (Proposition \texttt{prop:Treh}),
and $\mathcal{P}_s = 2.10\times10^{-9}$ from gravitational reheating.

\textbf{Thawing dark energy.}
The most immediately verifiable prediction is the thawing dark-energy relation
$w_a = -3(1+w_0)$ (Theorem \texttt{thm:thawing}), proved from the $k$-essence
equation of state in the slowly-rolling limit. This gives a one-parameter
prediction for the DESI DR2 $(w_0, w_a)$ plane. The DESI DR1 data are
consistent at $0.80\sigma$; DESI DR2 will confirm or falsify this at the
few-sigma level.

\subsection{Papers VIII--IX: Bell inequalities and the Born rule}

\textit{Prerequisites:} Bell inequalities; CHSH formalism; division algebras
$\HH$, $\CC$, $\OO$; Hopf fibration $S^3\to S^2$; Fourier analysis on $S^1$.

These two papers address the most conceptually striking aspect of the series:
the framework's relationship to quantum mechanics. The suppression law
$S(\theta) = \sin^4\theta/\vp^6$ is not just a Bell-violation weighting; its
$\sin^4\theta$ factor is the squared quaternionic norm $|q|^4$ for $q\in S^3$,
mapped to $S^2$ by the Hopf fibration.

The central results:
\begin{itemize}[noitemsep]
\item The condensate correlation $\xcond(\theta) = 2\cos\theta/(1+\cos^2\theta)$
  is the push-forward of $\cos\theta$ under the Hopf map (Paper~IX~\cite{Paper9} Theorem~3.2).
  The Tsirelson gap $2\sqrt{2} - 2.363$ decomposes exactly into a Fourier harmonic
  mismatch and a geometric projection error from the $\HH\to\CC$ averaging.
\item The Born rule $P=|\psi|^2$ emerges from the condensate amplitudes when
  complexified via the topological winding phase $\Phi = Q\omega$ (Theorem~5.2), proved as a mathematical theorem in Paper~IX~\cite{Paper9} and
  justified as a physical statement by the canonical quantisation of Paper~X~\cite{Paper10}.
\item The photon CHSH ceiling is exactly $\Tsirelson$ (Theorem~9.14),
  reproduced by the Stokes/Malus observable with $m=2$ Fourier mode.
  The lepton CHSH ceiling is $8/3 \approx 2.667$ (Proposition~8.4),
  a distinct prediction for loophole-free lepton Bell tests. Both in Paper~IX~\cite{Paper9}.
\end{itemize}

The Malus interaction Hamiltonian
$\langle\phi_n|H_{\mathrm{int}}|\alpha\rangle = g_{\mathrm{eff}}\cos(\phi_n-\alpha)$
with $g_{\mathrm{eff}} \propto \vp^{-6}$ is derived from the Noether current of the
$k$-essence action (Paper~IX~\cite{Paper9}. Theorem~9.8). Malus's law and the photon winding
number $m_{\mathrm{phot}}=2$ are theorems, not inputs.

\subsection{Paper X: Canonical quantisation}

\textit{Prerequisites:} Geometric quantisation (Kirillov--Kostant--Souriau);
Dirac constraint quantisation; spectral zeta functions; Macfarlane--Biedenharn
$q$-deformed oscillators; Bekenstein--Hawking entropy; Wald entropy formula.

Paper~X~\cite{Paper10} completes the series by constructing the full quantum theory. The
parent $k$-essence Lagrangian determines a symplectic structure $\Omega$ on the
field configuration space. After quotienting by Derrick dilations (the only
negative-Hessian direction, Theorem~4.2), the reduced space
admits prequantisation (Theorem~5.1) and geometric quantisation
(Theorem~5.3).

The resulting Hilbert space is $\mathcal{H} = L^2(\mathcal{C}_Q^{\mathrm{red}}, d\mu)$,
the spectrum is $E_n = \vp^{2n}E_0$ (Theorem~6.3), and the
$q$-deformation parameter is $q=\vp^2$ with $q$-numbers $[n]_{\vp^2} = F_{2n}$
(the even Fibonacci numbers, Theorem~9.4). Transition amplitudes
are exact Fibonacci ratios, making the theory exactly solvable.

The Bekenstein--Hawking entropy is derived in two independent steps: the Wald
entropy formula gives $S_{\mathrm{Wald}} = A_H/(4G_N)$ exactly at tree level
(Theorem~10.5), and the one-loop correction vanishes exactly
$\delta S_{1\text{-loop}} = 0$ because the spectral zeta function
$\zeta_{\mathcal{L}_*}(s) = (\vp^{2s}-1)^{-1}\zeta_{\mathrm{Riemann}}(s)$ gives
a temperature-independent effective action (Theorem~10.1).

\subsection{Paper XI: The $\varphi$-spiral and the periodic table}

\textit{Prerequisites:} Atomic units; Koopmans' theorem; Slater--Condon
rules; Hartree--Fock theory.

Paper~XI~\cite{Paper11} is a corollary of Papers~III and~IV rather than a new physical
framework. Its central observation: since $\aem$ and $m_e$ are both derived
from $\TCMB$ alone, the Rydberg energy $E_{\mathrm{Ry}} = \aem^2 m_e/2$
is a \emph{fixed, calculable tower level} $n(E_{\mathrm{Ry}}) = -12.102$
with no free parameters. The entire atomic energy scale is therefore a
prediction of the condensate, not an input.

Placing all first ionisation energies of the elements on the
$\vp$-spiral $r = e^{n\ln\vp}$ produces three structural observations:
(i) the visible electromagnetic spectrum ($1.77$--$3.1\,\mathrm{eV}$)
coincides with the pure tower level $\mathcal{L}_{\mathrm{cond}}\vp^{20}
= 1.799\,\mathrm{eV}$ to $1.7\%$, the same index as the electron mass;
(ii) hydrogen and oxygen are near-degenerate on the spiral
($\Delta n = 0.0015$, $0.15\%$), a geometric consequence of the Rydberg
being a fixed tower level;
(iii) Goldschmidt geochemical substitution pairs cluster within
$\Delta n < 0.004$ on the spiral.

The key new derivation (Propositions 5.1 and 5.3) proves $Z_{\mathrm{eff}}(\mathrm{O})=2$
in two levels: Level~1 from $Z=8$ and the Slater--Condon exchange algebra alone
(no numerical HF); Level~2 from the WZW Chern--Simons correction of Paper~IV~\cite{Paper4}
giving the $0.087\%$ shift to $Z_{\mathrm{eff}}=2.00087$.

\subsection{Paper XII: Profile normalisation and the one-parameter chain}

\textit{Prerequisites:} Papers~I--III (gradient energy integrals $J_{2a}$,
$J_4$; virial self-consistency $V(b^*)=\vp$; Linking Scale ratio
$J_4/J_{2a}=2^{4/3}/\vp^5$); thermodynamics of a field in a radiation bath.

Paper~XII~\cite{Paper12} closes the last logical gap in the series by proving the profile
normalisation conjecture $\vp^9\sqrt{J_{2a}J_4}=10$, the step needed to
make $m_e$ derivable from $\TCMB$ alone.
The proof is short and uses no numerical input.

\textbf{The key identity.}
The condensate scale satisfies $\rho_{\mathrm{CMB}} = 10\,\Lambda_{\mathrm{cond}}^4$
exactly: from $\Lambda_{\mathrm{cond}} = \TCMB(\pi^2/150)^{1/4}$ and
$\rho_{\mathrm{CMB}} = (\pi^2/15)\TCMB^4$, the ratio is $150/15=10$,
a pure arithmetic identity.
The integer $10 = \mathrm{ord}(q_{2I})$ is therefore literally the CMB
energy density in natural units.

\textbf{The equivalence.}
The paper proves that the conjecture is \emph{equivalent} to the thermodynamic
condition $E_{\mathrm{fb}}/V_{\mathrm{torus}} = \rho_{\mathrm{CMB}}$:
the condensate's gradient energy density equals the ambient CMB energy density.
Using the proved ratio $J_4/J_{2a}=2^{4/3}/\vp^5$, both sides reduce to the
geometric constraint $R_0 C^{*2} = 5/(\vp^{23}\pi^2)$.
The equilibrium condition holds because the condensate torus adjusts its size
until its energy density matches the CMB background --- the standard
thermodynamic condition for a field embedded in a radiation bath.
Combined with $\rho_{\mathrm{CMB}}=10\,\Lambda_{\mathrm{cond}}^4$, this forces
the conjecture.

\textbf{The one-parameter chain.}
\[
  \TCMB
  \;\xrightarrow{\Lambda_{\mathrm{cond}}}\;
  \rho_{\mathrm{CMB}}=10
  \;\xrightarrow{E_{\mathrm{fb}}/V=\rho_{\mathrm{CMB}}}\;
  \vp^9\sqrt{J_{2a}J_4}=10
  \;\xrightarrow{\text{Paper~I~\cite{Paper1}}}\;
  m_e = \Lambda_{\mathrm{cond}}\,\vp^{20}\,e^{4\pi-3/400}.
\]

\textbf{Additional results.}
(i)~Pentagon identities $\vp=2\cos(\pi/5)$, $\sin(\pi/5)=\frac{1}{2}\sqrt{\mu^*}$,
$\sin(2\pi/5)=\frac{\vp}{2}\sqrt{\mu^*}$ (where $\mu^*=3-\vp$).
(ii)~Quadratically-convergent Newton--icosahedral series
$q^*=\pi/5+d_1+d_2+\cdots$, $b^*=\tan^2(q^*)/8$, with first correction
$d_1=\frac{\vp\sqrt{\mu^*}}{12}\cdot\frac{8\sqrt{\mu^*}-3\pi}{5\vp\sqrt{\mu^*}/4-\pi}$
--- an alternating $(\vp,\pi)$-series structurally parallel to the $\alpha^{-1}$
formula of Paper~IV~\cite{Paper4}.
(iii)~Exact toroidal integral
$I_{\mathrm{tor}}(x)=\frac{2}{15}\,{}_2F_1(2,2;\frac{7}{2};-x^2)$,
with the exact evaluation $I_{\mathrm{tor}}(1)=\pi-3$ at equal radii $x=1$.

\subsection{Paper XIII: Atomic quantum mechanics from the condensate}

\textit{Prerequisites:} Paper~X~\cite{Paper10} (canonical quantisation, Hilbert space);
Paper~VII~\cite{Paper7} (chameleon screening); Paper~I~\cite{Paper1} (Hopf charge $Q=10$, $\twoI$ symmetry);
non-relativistic quantum mechanics (Schr\"odinger equation, angular momentum,
Hartree--Fock theory).

Paper~XIII~\cite{Paper13} proves that standard atomic and molecular quantum mechanics follows
from the condensate quantum theory (Paper~X) without additional postulates.
It contains ten main results in three groups.

\textbf{Group 1: The Schr\"odinger equation.}
The one-particle sector of the Paper~X~\cite{Paper10} Hilbert space, restricted to the
translation fibration of the condensate, gives the momentum operator
$\hat{p} = -i\hbar\nabla$ as the Noether current (Lemma \texttt{lem:momentum}).
Born--Oppenheimer separation of the slow (nuclear) and fast (condensate)
degrees of freedom yields the time-independent Schr\"odinger equation
$H\psi = E\psi$ with $H = -\hbar^2\nabla^2/(2m_e) + V(r)$ as a mathematical
theorem (Theorem~7.4). The non-adiabatic coupling
vanishes \emph{exactly} (Theorem~7.7), not as an
approximation --- a consequence of the condensate's spectral gap $\vp^2-1$.

\textbf{Group 2: Spin and statistics.}
The electron is a $Q=2$ (Paper~I~\cite{Paper1} gradient-knot. Its winding structure
under $\twoI$ gives $\pi_1(\mathrm{SO}(3)) = \ZZ_2$ holonomy, proving
spin-$\tfrac{1}{2}$ as a topological theorem (Theorem~8.1).
Fermi statistics and the Pauli exclusion principle follow as corollaries.
The complete spin tower $J = Q/4$ (Theorem~8.6) assigns
$J=0$ to $Q=0$ (vacuum), $J=1/2$ to $Q=2$ (electron), $J=1$ to $Q=4$ (photon),
$J=3/2$ to $Q=6$ (gravitino), $J=2$ to $Q=8$ (graviton).

\textbf{Group 3: Atomic structure.}
Chameleon screening (Paper~VII~\cite{Paper7}) at atomic densities exponentially suppresses
the condensate correction to $O(e^{-\rho_{\rm atom}/\rho_\infty})$, proving
that the Coulomb potential and spherical harmonics emerge from first principles
(Paper~XII~\cite{Paper13}~Theorem~4.1). The result extends to
multi-electron atoms: the Hartree--Fock equations are exact within the
condensate framework (Paper~XII~\cite{Paper13}, Corollary~4.6).

\textbf{Testable prediction.}
The free-electron mass anisotropy in the CMB rest frame
(Proposition~6.1):
\begin{equation}
  \frac{\Delta m}{m_e} = \frac{1}{\vp^8} \approx 2.13\%,
\end{equation}
measurable via Penning trap cyclotron frequency comparison at
$\delta\omega/\omega \sim 10^{-4}$.

\subsection{Paper XIV: Thick-torus toroidal correction and electroweak closure}

\textit{Prerequisites:} Papers~I--III (gradient energy integrals $J_{2a}$,
$J_4$; modified BPS equation; WZW normalization);
Paper~IV~\cite{Paper4} (Yukawa coupling formula, $\vEW$ prediction);
toroidal geometry; azimuthal averaging.

Paper~XIV~\cite{Paper14} closes the last numerical gap in the series: the $0.69\%$
residual between the predicted and measured electroweak VEV.
The thin-torus approximation used in Paper~IV~\cite{Paper4} gives
$\vEW^{\rm pred} = 247.9\,\mathrm{GeV}$ ($+0.69\%$ above the measured
$246.2\,\mathrm{GeV}$). Paper~XIV~\cite{Paper14} computes the $O(R_0^{-2})$ thick-torus
correction exactly.

\textbf{The exact azimuthal reduction.}
All $\chi$-integrals on the condensate torus (major radius $R_0=3$,
minor radius $\rho$) are evaluated in closed form via three identities:
$\langle r\rangle = R_0$, $\langle 1/r\rangle = 1/\sqrt{R_0^2-\rho^2}$,
$\langle 1/r^2\rangle = R_0/(R_0^2-\rho^2)^{3/2}$ (Proposition~4.1).

\textbf{The closed-form toroidal correction.}
The 3D Faddeev--Niemi integrals decompose as
(Theorem~4.3):
\begin{align}
  J_4^{3D} &= R_0\,I_4 + (\vp^8+1)\,I_{2a}^{\rm tor}, \\
  J_{2a}^{3D} &= R_0\,I_{2a} + (\vp^8+1)\,I_4^{\rm tor},
\end{align}
where $Q_t^2 = \vp^8+1 \approx 47.979$ is the toroidal winding parameter
and $I_n^{\rm tor}$ are 1D integrals with weight
$\rho/\sqrt{R_0^2-\rho^2}$. The cross-coupling structure (azimuthal
winding feeds $J_4$; Hopf density feeds $J_{2a}$, swapped from the
1D assignment) resolves the factor-of-4.2 discrepancy between thin-torus
and full 3D computations. At $C^*=3.4318$, $R_0=3$:
$J_4^{3D}/J_{2a}^{3D} = 0.22717$ versus WZW target $0.22721$ ($0.018\%$).

\textbf{Yukawa correction and closure.}
The Paper~I~\cite{Paper1} WZW formula $d\ln y_e/dC^* = 1/(2k) = 1/6$ combined with
$\delta C^* = C^{*,3D} - C^{*,{\rm thin}} = 3.432 - 3.392 = 0.040$
gives (Corollary~4.7):
\begin{equation}
  \frac{\delta y_e}{y_e} = \exp\!\left[\frac{\delta C^*}{2k}\right] - 1
  \approx +0.68\%,
\end{equation}
which reduces $\delta\vEW/\vEW$ from $+0.69\%$ to $0.04\%$ and
$\delta\Lambda_{\mathrm{QCD}}/\Lambda_{\mathrm{QCD}}$ from $+0.97\%$
to $0.04\%$. The toroidal parameter $\vp^8+1$ equals the Paper~XIII~\cite{Paper13}
tangential mass numerator, connecting the thick-torus geometry to the
free-electron mass anisotropy.

\textbf{The Higgs boson mass.}
The condensate's $Q=0$, $J=0$ sector (Paper~XIII spin tower $J=Q/4$) is the scalar
Higgs sector. The condensate fills space as a scalar background whose VEV generates
fermion masses; the Higgs boson is the radial breathing mode of the condensate around
this equilibrium, not a topological excitation.
The quartic self-coupling is (Theorem~6.11):
\begin{equation}
  \lambda = \frac{k+2}{2}\,r_{\mathrm{WZW}}^2
  = \frac{5\cdot 2^{5/3}}{\vp^{10}} = 0.12907,
\end{equation}
giving Higgs mass $\mH = \sqrt{k+2}\,r_{\mathrm{WZW}}\,\vEW
= \sqrt{5}\,(2^{4/3}/\vp^5)\,\vEW = 125.10\,\mathrm{GeV}$,
which is $0.13\sigma$ from the ATLAS combined Run~1+2 measurement
$125.11\pm0.11\,\mathrm{GeV}$ (residual $-0.01\%$).
The derivation involves two layers.
The \emph{profile layer}: the ratio $r_{\mathrm{WZW}} = J_4/J_{2a} = 2^{4/3}/\vp^5$
is the same stiffness ratio that governs the Yukawa coupling and the toroidal
correction; it determines $\lambda \propto r_{\mathrm{WZW}}^2$.
The \emph{WZW layer}: the normalization factor $(k+2)/2 = 5/2$ equals the
pentagon Parseval sum $\sum_{n=1}^{k+1}\sin^2(n\pi/(k+2)) = (k+2)/2$
(a consequence of the $(k+2)$-th roots of unity summing to zero), proved to be
the correct torus four-point normalization via the KZB $b$-cycle monodromy
(Falceto--Gaw\c{e}dzki 1996): on the torus, periodic conformal block boundary
conditions replace the plane OPE weights $(C^j)^2 d_j$ with the modular
weights $d_j^2\sin^2(\pi/(k+2))$, which sum to $(k+2)/2$ by the Parseval
identity.
The plane result $(3-\vp)r_{\mathrm{WZW}}^2$ gives $\mH \approx 93\,\mathrm{GeV}$;
the torus enhancement $\sqrt{5}\vp/2 \approx 1.809$ lifts this to $125.10\,\mathrm{GeV}$.

\subsection{Paper XV: The $Q_H=3$ trefoil sector and quark charges}

Paper~XV extends the framework from the $Q_H=2$ lepton sector (torus, $\twoI$,
$\mathrm{SU}(2)_3$) to the $Q_H=3$ quark sector (trefoil knot $T_{2,3}$,
binary tetrahedral group $2T$, $\mathrm{SU}(3)_1$).

\textit{QFT/GR prerequisites:} WZW models and conformal field theory at level~1;
Verlinde formula and modular S-matrix; McKay correspondence for finite
subgroups of $\mathrm{SU}(2)$; basic knot theory (trefoil, writhe, Hopf invariant);
fusion categories at level $k=3$.

\textbf{The McKay chain.}
The binary tetrahedral group $2T\subset\twoI$ ($|2T|=24$, index~5)
has McKay graph $\widehat{E}_6$ (7~nodes), giving $(E_6)_1$ as the
McKay-direct WZW model ($c=6$).
The conformal embedding $(E_6)_1\supset\mathrm{SU}(3)_1^{\times 3}$
via the trinification $E_6\supset\mathrm{SU}(3)_C\times\mathrm{SU}(3)_L\times\mathrm{SU}(3)_R$
provides an exact central-charge match $c=6=3\times 2$.
This is mathematically stronger than the pentagon identification
$k+2=5$ used for $Q_H=2$: the latter has no conformal embedding
from $(E_8)_1$ to $\mathrm{SU}(2)_3$ (different central charges $c=8$ vs.\ $9/5$).

\textbf{Quark charges.}
The trinification charge generator $Q=T_{3L}+T_{8L}/\!\sqrt{3}$
is uniquely fixed by $\mathrm{SU}(3)_L$ group theory (tracelessness plus
$\mathbb{Z}_3$ centre compatibility; no SM hypercharge matching is used).
The charge-generator direction $\mu=2\omega_1$ and the $\mathrm{SU}(3)_1$
conformal weight $h_{(1,0)}=(\omega_1,\omega_1)_{A_2}/2=1/3$
give both quark charges from the single number
$(\omega_1,\omega_1)_{A_2}=2/3$:
$|Q_d|=h_{(1,0)}=1/3$, $|Q_u|=2h_{(1,0)}=2/3$.
This holds only at $N=3$: the unique $N$ with
$h_{\mathrm{fund}}(\mathrm{SU}(N)_1)=1/N$.

\textbf{The Bogomolny parameter $\lambda_3=\vp^6$: proved.}
The $j=3/2$ simple current of $\mathrm{SU}(2)_3$ has quantum dimension~$1$,
self-fusion $\tfrac{3}{2}\otimes\tfrac{3}{2}=0$ ($N_{\mathrm{fus}}=1$),
and the sector assignment $Q_H=3\leftrightarrow j=3/2$ is proved
by three independent arguments:
(i)~triple-fusion multiplicity plus $\pi_3(S^2)=\mathbb{Z}$
topological uniqueness;
(ii)~Verlinde uniqueness ($|S_{j,1/2}/S_{j,0}|=\vp$ only at
$j=0$ and $j=3/2$);
(iii)~$\mathbb{Z}_2$ charge parity pre-filtering.
The LSC formula then gives $s_{\mathrm{opt}}^{2k}=1$,
hence $\lambda_3=\vp^6$ in three lines.
All previously conditional results ($r_3=2/\vp^5$,
$\vp^9\sqrt{J_{2a}J_4}=3$, $V_3=9/(5\vp^{17})$, $C_3^*=2.5062$)
are now unconditional.

\textbf{The second-order curvature correction.}
The leading-order trefoil formula $r_3 = 3/(4C_3^{*2})+\langle\kappa^2\rangle/2$
gives $C_3^*=2.497$. The $O(\langle\kappa^4\rangle)$ correction is computed
exactly via Beta-function integration of the BS profile in normalised tube
coordinates: $\delta r_3^{(\kappa^4)} = 9\langle\kappa^4\rangle/(4 C_3^{*4})
= 0.000938$ for $T_{2,3}$, shifting the saddle parameter to
$C_3^*=2.5062$ ($0.392\%$ shift from the leading-order value). This is the
trefoil-sector analogue of Paper~XIV's thick-torus correction, and likewise
leaves $\lambda_3$ unaffected: the Bogomolny parameter is fixed independently
by the sector-assignment proof, while the curvature correction refines only
the geometric saddle parameter $C_3^*$.

Three exact WRT results provide independent corroboration:
the framing phase $\theta_{3/2}^3=i$ (primitive 4th root),
the quantum dimension $\dim_q(3/2;3)=1$ (string tension minimised),
and the pure-unit-modulus WRT contribution ($|$WRT$|=1$, no
dimensional renormalisation under trefoil surgery).

\subsection{Paper XVI: The quark generation-mass search}

\textit{Prerequisites:} Everything required for Paper~XV, plus one-loop
multi-threshold QCD running; basic Beta-function/gradient-flow numerics.

Paper~XV fixes the qualitative quark structure but leaves the
generation-mass hierarchy open: the inter-generation WZW T-matrix phases
give only an $O(1)$ correction, so the dominant six-order-of-magnitude
hierarchy from $m_u$ to $m_t$ needs a mechanism not supplied by Paper~XV.
Paper~XVI is a systematic negative-results survey: six candidate
mechanisms are tested and each is ruled out for a specific, stated
structural reason (wrong order of magnitude, a symmetry that fails to
supply an independent generation label, a level-rigidity obstruction,
an asymmetric coset splitting, a residual that is not a smooth function
of generation index, and so on), rather than a vague failure to match
data.

\textbf{The Y-junction candidate.} A seventh, qualitatively distinct
mechanism, motivated by colour confinement, is examined rather than
ruled out: the trefoil's Y-junction string-tension energy
$\sigma\cdot3R_0$, present only for $Q_H=3$ and absent for the
unconfined $Q_H=2$ torus. The chameleon-sector scalar field and the
Hopf profile function are shown to be the same field under an explicit
map, removing one possible obstruction to the idea, but the Y-junction
arms lie off the trefoil's own tube, so the string tension cannot be
read off existing tube integrals. Both the string-tension energy and
the junction energy are nevertheless computed --- the former
analytically, the latter numerically via gradient flow in the junction
domain --- and the junction energy itself is ruled out as a source of
generation-dependence ($E_\cup/E_{\mathrm{baryon}}\sim10^{-7}$).

\textbf{QCD running of the colour shift.} Paper~V's proved colour-level
shift $\delta n^{(C)}=\tfrac12$ is run, for the first time with a
proper one-loop multi-threshold treatment, from the condensate UV scale
to the PDG comparison scale. The strange-quark prediction matches
measurement to $0.09\%$; the other five generation predictions remain
off by factors of $1.7$ to $49$, in a systematic, generation-growing,
always-undershooting pattern --- consistent with, but not yet derived
from, a missing confinement-related energy contribution.

A QFT reader should treat this paper as a record of a genuinely open
problem, documented so that later work does not re-walk the same ruled-out
paths; it does not itself close the quark-mass hierarchy.

\subsection{Paper XVII: The $Q_H=1$ sector and the neutrino}

\textit{Prerequisites:} Everything required for Papers~XV--XVI; A-type
ADE/McKay classification; PMNS mixing formalism.

With $Q_H=3$ (quarks/baryons) and $Q_H=2$ (charged leptons) identified in
Papers~XV--XVI, Paper~XVII completes the sector inventory by
characterising $Q_H=1$. Four independent structural results identify it
as the neutrino sector:

\begin{itemize}[noitemsep]
\item \textbf{Topology.} The $Q_H=1$ preimage of a regular value on $S^2$
  is a single unknotted circle with writhe zero, in contrast to the
  $Q_H=2$ torus and $Q_H=3$ trefoil.
\item \textbf{Group theory.} $Q_H=1$ is organised by the cyclic group
  $\Ztwo$ via the A-type McKay correspondence
  $\Ztwo\leftrightarrow\widehat{A}_1$, rather than the E-type groups
  governing the higher sectors. Since $\Ztwo$ and the colour $\Zthree$
  are coprime, colour exclusion for $Q_H=1$ follows immediately,
  strengthening the colour-exclusion theorem of Paper~XVI.
\item \textbf{WZW model.} The associated model is $\mathrm{SU}(2)_1$
  ($c=1$, two primaries), giving $Q_{\mathrm{group}}=6$ --- the smallest
  in the sector hierarchy $\{3,6,10\}$ for $\{Q_H=3,Q_H=1,Q_H=2\}$.
\item \textbf{Quantum dimension.} At $k=3$ the $j=\tfrac12$ primary has
  $\dim_q=\vp$, matching the $Q_H=2$ sector and confirming $Q_H=1$ is
  unconfined (a confined sector requires $\dim_q=1$).
\end{itemize}

\textbf{Mass prediction.} The parameter-free formula
$m_{\nu_1} = \TCMB(\pi^2/90)^{1/4}\cdot\vp^{12}\cdot e^{-5/144}
\approx 42\,\mathrm{meV}$ follows from the same CMB-matching method as
the electron mass (Paper~XII). It sits within the final target
sensitivity of the Project~8 experiment ($m_\beta\geq40\,\mathrm{meV}/c^2$
at $90\%$~CL); a non-detection at that sensitivity would falsify the
$Q_H=1$ identification. The tower-level bound $n_\nu\geq42.02$ resolves
a previous unwarranted integer-rounding artefact via Paper~VI's
irrationality theorem for tower levels. The Wess--Zumino descent forces
$k=1$ for the $Q_H=1$ condensate, and the resulting self-conjugacy of
the $j=\tfrac12$ primary predicts Majorana neutrinos.

\textbf{PMNS matrix.} The PMNS mixing matrix is derived exactly as a
permutation, and $Q_{\mathrm{eff}}$ is derived from Brieskorn geometry.
The remaining open problem is the $q$-series resummation converting
branching functions into physical mixing angles beyond the coset
structure.

%═══════════════════════════════════════════════════════════════════════════════
\section{Does this framework contradict established physics?}
\label{RG:sec:consistency}
%═══════════════════════════════════════════════════════════════════════════════

\subsection{The short answer}

No. The framework as constructed \emph{reproduces} established observations,
and does not replace them. Every SM observation it addresses is reproduced to at least
the accuracy quoted in Papers~I--XVII. No observed quantity is contradicted.

The framework's novelty is not contradiction but compression: it reduces the
Standard Model's $\sim30$ free parameters to one measured input.

\subsection{Particle physics: what is derived vs measured}

The table below gives the main SM-sector derivations and their residuals.

\begin{center}
\begin{longtblr}[
  caption = {},
  label = {tab:particle_physics},
] {
  colspec = {X[2,p] X[3.5,p] X[1,c] X[1,c]}, % Proportional widths
  rowhead = 1, % Repeat header on each page
}
\toprule
Observable & Source in series & Derived & Residual \\
\midrule
+$\aem^{-1}$ (three-term) & P4, Thm~5.1 & 137.036\,49 & $0.0004\%$ \\
+$\aem^{-1}$ (four-term) & P4, Thm~5.3 & 137.036\,00 & $5\times10^{-7}\%$ \\
$m_e$ & P12, profile norm & 0.511 MeV & $0.013\%$ \\
$m_\mu/m_e$ & P6, WZW T-phase $T_1,T_2$ & 206.77 & $0.22\%$ \\
$m_\tau/m_e$ & P6, WZW T-phase $T_3$ & 3477 & $0.34\%$ \\
Koide sum rule & P1, $\twoI$ characters & exact & --- \\
$\alpha_s(M_Z)$ & P5, semi-Dirac $d_{\mathrm{eff}}$ & 0.1185 & $2.4\%$ \\
$m_u/m_d$ & P5, $E_8$ Coxeter exponents & within 1-loop & $<1\%$ \\
CP phase $\delta_{\mathrm{CP}}$ & P5, holonomy & $71.2°$ & --- \\
$|V_{cb}|$ & P5, WZW fusion & $0.04068$ & $0.4\%$ \\
$n_s$ (CMB) & P7, Starobinsky branch & 0.9654 & $0.12\sigma$ \\
$r$ (tensor-to-scalar) & P7, reheating & 0.0036 & --- \\
$w_a$ (dark energy) & P7, thawing & $-3(1+w_0)$ & $0.80\sigma$ (DR1) \\
$Z_{\mathrm{eff}}(\mathrm{O})$ & P11, Koopmans & $2.00087$ & $0.003\%$ \\
Rydberg tower index & P11, corollary & $n(E_{\mathrm{Ry}})=-12.102$ & exact \\
Schr\"odinger eqn & P13, Thm & derived & --- \\
Spin-$\frac{1}{2}$ & P13, Thm & $J=Q/4$ & --- \\
$\Delta m/m_e$ (anisotropy) & P13, Prop & $1/\vp^8 \approx 2.13\%$ & untested \\
$\vEW$ & P14, Cor & $246.3\,\mathrm{GeV}$ & $0.04\%$ \\
$\Lambda_{\mathrm{QCD}}$ & P14, tower & $208.0\,\mathrm{MeV}$ & $0.04\%$ \\
$\mH$ ($\vEW$ measured) & P14, Thm & $125.10\,\mathrm{GeV}$ & $0.13\sigma$ \\
$\mH$ (1-input) & P14, Thm & $125.11\,\mathrm{GeV}$ & $0.08\sigma$ \\
$|Q_d|$ (down-quark charge) & P15, trinification & $1/3$ & exact \\
$|Q_u|$ (up-quark charge) & P15, trinification & $2/3$ & exact \\
$m_s$ (strange quark) & P16, colour-shift QCD running & matched & $0.09\%$ \\
$m_\nu$ tower bound & P17, irrationality theorem & $n_\nu\geq42.02$ & margin $0.37$ \\
$m_{\nu_1}$ (neutrino mass) & P17, CMB matching & $42\,\mathrm{meV}$ & untested (Project~8) \\
\bottomrule
\end{longtblr}
\end{center}

The residuals are not zero. The framework states this explicitly: at $k=3$
the formula gives $\aem^{-1} = 137.036$, not $137.03599908\ldots$; the
$5\times10^{-7}\%$ residual is acknowledged as a beyond-one-loop WZW correction not
yet computed. Similarly for lepton masses. These are genuine open questions,
not mismatches: the series never claims exact agreement beyond the order at
which the calculation is performed.

\subsection{What the framework does not predict (and why this is consistent)}

\textbf{Neutrino masses.} The condensate framework predicts a normal hierarchy
with lightest neutrino mass $\sim3\times10^{-3}$\,eV (Paper~VII~\cite{Paper7}, Proposition
12.1), but the absolute mass scale requires input from
cosmological data not yet available at the sensitivity needed.

\textbf{Dark matter direct detection.} The condensate acts as dark matter via
clustering of condensate perturbations (sound speed $c_s = 1/\vp$ at the
attractor density), not through a new particle. It therefore predicts \emph{null
results} in all dark-matter direct-detection experiments, consistent with all
existing constraints.

\textbf{Deviation from GR at laboratory scales.} The Vainshtein mechanism
(Paper~VII~\cite{Paper7}, Theorem~6.9) suppresses scalar fifth-force
corrections to Newtonian gravity at Solar System scales to below Cassini bounds.
At laboratory scales the condensate is effectively GR.

\textbf{New particles at the LHC.} The framework does not predict any new
resonance below the condensate UV cutoff $\Lambda_{\mathrm{UV}} \approx
0.0557\,M_{\mathrm{Pl}} \approx 10^{17}$\,GeV. The LHC energy scale is many
orders of magnitude below this threshold.

\subsection{The genuine tensions and open questions}

Being precise about what is unresolved matters for a fair evaluation.

\textbf{The $5\times10^{-7}\%$ residual in $\aem$.} The three-term formula is a
one-loop WZW result. The four-term formula (Paper~IV~\cite{Paper4}, Theorem~5.3)
reduces the residual from $0.0004\%$ to $5\times10^{-7}\%$
($\alpha^{-1} = 137.036\,00$ vs measured $137.036\,00$).
The fourth term $-1/(36\pi\vp^6) = -T_3/(4\pi)$ is the Chern--Simons
correction to the QED vacuum polarisation inside the condensate background:
the SU(2)$_3$ magnetic flux threads the QED loop, suppressing it by
the CS coupling $1/(4\pi)$ per unit WZW level.
The series $T_4/T_3 \approx 1/(4\pi) \approx 0.08$ converges geometrically;
Paper~IV estimates a further term at this rate would be
$\sim6\times10^{-6}$ in $\alpha^{-1}$ (Paper~IV, Remark 5.4), for
illustration of convergence only --- the WZW series itself terminates
at four terms (two-loop WZW vanishes exactly). The current four-term
residual, $6.9\times10^{-7}$, is well within experimental reach
(current precision $\sim1$--$3\times10^{-8}$, e.g.\ Morel et al.\ 2020);
Paper~IV attributes it to ordinary sub-leading QED physics outside the
condensate framework, and separately treats it as an open question with
an as-yet-unverified candidate closed form.

\textbf{The $\sigma_{\mathrm{xc}}$ derivation (equation 25) in Paper~XI~\cite{Paper11}.} The exchange shielding constant
$\sigma_{\mathrm{xc}} = 1.55$ for the oxygen 2p orbital is confirmed by the
Clementi--Raimondi Hartree--Fock table but is not derived from first principles
within the framework. The formula is correct; its derivation is incomplete. This
is acknowledged in Paper~XI~\cite{Paper11}.


\textbf{The lepton CHSH ceiling.} The prediction $\CHSH_{\mathrm{lepton}} = 8/3$
for loophole-free lepton Bell tests has not yet been tested. No existing Bell
experiment uses lepton spin in the geometry required.

\textbf{The BAO peak shift.} The prediction $k_{\mathrm{BAO}}^{\mathrm{cond}}
= \vp \times k_{\mathrm{BAO}}^{\mathrm{CDM}} \approx 0.0346\,\mathrm{Mpc}^{-1}$
distinguishes the condensate from standard CDM. It is falsifiable with DESI DR2
and Euclid.

\textbf{The WZW bosonization in 3+1D.}
Paper~III~\cite{Paper3}, Theorem~10.3
proves the 3+1D extension of the bosonization equivalence with explicit
error bounds. The effective $\mathrm{SU}(2)_3$ WZW action arises from the
zero mode of the Kaluza--Klein reduction at energies $E\ll r_{\mathrm{core}}^{-1}$,
with corrections $O(R_0^{-2})$ matching the rate established independently
by Proposition~8.3.
The Wess--Zumino term descends correctly via the Zumino--Faddeev descent
equations, the Hopf charge $Q=10$ is preserved by a homotopy argument,
and the $\mathrm{SU}(2)$ 't~Hooft anomaly matches at level $k=3$.
The remaining open item (noted in Remark~10.5)
is the convergence of the KK mode expansion for non-static configurations;
this affects the $O(R_0^{-2})\approx0.22\%$ subleading corrections
but not the leading-order SM predictions.

%═══════════════════════════════════════════════════════════════════════════════
\section{Technical prerequisites: a compact reading list}
%═══════════════════════════════════════════════════════════════════════════════

A reader with standard QFT/GR training will be comfortable with the following
material and can enter the series at the level stated.

\begin{center}
\begin{longtblr}[
  caption = {},
  label = {tab:reading_list},
] {
  colspec = {X[3.8,p] X[5.5,p] X[3.5,p]}, % Proportional widths
  rowhead = 1, % Repeat header on each page
}
%\begin{tabular}{p{3.8cm}p{5.5cm}p{3.5cm}}
\toprule
Topic & Needed for & Standard reference \\
\midrule
Faddeev--Niemi functional & Papers I, II & Faddeev--Niemi (1997) \\
WZW models and $\mathrm{SU}(N)_k$ & Papers I, III, IV, VI & Di~Francesco et~al.\ (CFT book) \\
Verlinde formula & Papers I, III, IV, V & Verlinde (1988) \\
Derrick--Pohozaev identity & Papers II, VI & Derrick (1964) \\
Non-abelian bosonization & Paper III & Witten (1984) \\
Modular T-matrix / T-phases & Papers IV, VI & Di~Francesco et~al.\ \\
McKay correspondence & Papers I, FC & McKay (1980) \\
$k$-essence / scalar-tensor & Papers VI, VII & Armendariz-Picon et~al.\ (2001) \\
Chameleon / Vainshtein & Paper VII & Khoury--Weltman (2004) \\
Starobinsky inflation & Paper VII & Starobinsky (1980) \\
Bell inequalities / CHSH & Papers VIII, IX & Bell (1964); CHSH (1969) \\
Division algebras ($\HH$, $\CC$, $\OO$) & Papers IX, FC & Baez (2002) \\
Geometric quantisation & Paper X & Woodhouse (1992) \\
Spectral zeta functions & Paper X & Berger--Gauduchon--Mazet \\
Wald entropy & Paper X & Wald (1993) \\
HF theory / Slater rules & Paper XI & Condon--Shortley (1935) \\
Goldschmidt geochemistry & Paper XI & Goldschmidt (1926) \\
Non-rel.\ quantum mechanics & Paper XIII & Griffiths (2018) \\
Born--Oppenheimer approx.\ & Paper XIII & Born--Oppenheimer (1927) \\
Toroidal geometry & Paper XIV & do Carmo (1976) \\
KZB equation on tori & Paper XIV & Falceto--Gaw\c{e}dzki (1996) \\
Trefoil knot theory & Paper~XV & Rolfsen (1976); Adams (1994) \\
Fusion categories / modular data & Paper~XV & Kitaev (2006); Di~Francesco et~al.\ \\
Conformal embeddings & Paper~XV & Di~Francesco et~al.\ Ch.~18; Bais--Bouwknegt (1987) \\
WRT knot invariants & Paper~XV & Witten (1989); Reshetikhin--Turaev (1991) \\
One-loop multi-threshold QCD running & Paper~XVI & Peskin--Schroeder Ch.~17 \\
ADE / McKay classification (A-series) & Paper~XVII & McKay (1980); Kac (1990) \\
PMNS matrix / neutrino mixing & Paper~XVII & Pontecorvo (1957); Maki--Nakagawa--Sakata (1962) \\
\bottomrule
\end{longtblr}
\end{center}

A reader who is unfamiliar with WZW models but strong in standard QFT should
read Di~Francesco et~al.\ Chapters~14--16 (the $\mathrm{SU}(2)_k$ WZW model,
the Kac--Moody algebra, the Verlinde formula, and the modular S and T matrices)
before Papers~III--IV. This takes roughly two days and unlocks $\sim$70\% of the
series.

A reader new to Paper~XV who is comfortable with Papers~III and~IV needs additionally:
Di~Francesco et~al.\ Chapter~18 (conformal embeddings, $E_6\supset\mathrm{SU}(3)^3$),
and the basic theory of fusion categories at $k=3$ (Kitaev 2006, Section~12).
The WRT invariant computation requires only the Verlinde S-matrix already used
in Paper~III; no additional background in knot theory is strictly needed to
follow the proofs, which reduce to combinatorial statements about fusion multiplicities.

%═══════════════════════════════════════════════════════════════════════════════
\section{Questions a QFT reader will ask: direct answers}
%═══════════════════════════════════════════════════════════════════════════════

\textbf{Is the condensate a physical field or a mathematical convenience?}
It is treated as physical: a topological soliton configuration of the vacuum
scalar sector. Paper~X~\cite{Paper10} constructs its quantum theory rigorously, with a Hilbert
space and a self-adjoint Hamiltonian. Whether the physical identification is
correct is an experimental question, not a mathematical one.

\textbf{Does it break Lorentz invariance?}
No. The condensate has a preferred direction only locally (the icosahedral
orientations of $\twoI$), and these average to an isotropic distribution
over macroscopic scales. The action~\eqref{RG:eq:action} is Lorentz-covariant
in $3+1$D; the density-feedback structure breaks Lorentz invariance only
in the sense that any non-trivial vacuum does.

\textbf{Is the bosonization in Paper~III~\cite{Paper3} rigorous?}
The Witten bosonization equivalence $\mathrm{SU}(N)_k\,\mathrm{WZW}
\leftrightarrow N$ free Dirac fermions is a proved theorem for $1+1$D QFT.
Paper~III applies it at the 2D cross-section of the condensate (Theorem
~10.1), proved rigorously via the LSC and Pohozaev identity.
The 3+1D extension uses dimensional reduction and is not proved with the same
rigour. A critical reader should treat the SM-sector predictions as conditional
on this extension being valid.

\textbf{Why does $\vp$ appear everywhere?}
The Foundational Companion answers this cleanly: all appearances of $\vp$
are consequences of $k+2=5$. This is not a coincidence or a design choice;
it is a theorem. The pentagon diagonal-to-side ratio is $\vp$ precisely
because $5$ is the number of sides, and $k=3$ forces $k+2=5$ through the
three routes described in Section~3.1.

\textbf{Are the SM numerical predictions genuine predictions or fits?}
The framework takes one inputs ($\TCMB$) and derives the rest.
There are no fitted parameters. However, the framework does make choices
(using the Witten bosonization dictionary, using the $E_8$ Coxeter assignment
for quarks, using the Starobinsky inflationary potential) that could in
principle be made differently. The correct way to think about this is:
the framework is a specific theoretical construction that makes definite,
falsifiable predictions. Its predictions agree with observation to the
accuracy computed. Whether a different construction with the same inputs
could also agree is a separate question.

\textbf{What would falsify it?}
The framework is falsifiable at multiple scales simultaneously:
\begin{enumerate}[noitemsep]
\item $w_a \neq -3(1+w_0)$ at DESI DR2 or Euclid at $>3\sigma$.
\item BAO peak position not shifted by factor $\vp$ at DESI DR2.
\item Loophole-free lepton Bell test giving $\CHSH > 8/3$ at $>5\sigma$.
\item CMB-S4 tensor-to-scalar ratio $r \gg 0.0036$.
\item A new SM-sector resonance at energies accessible to current experiments.
\item Free-electron mass anisotropy $\Delta m/m_e \neq 1/\vp^8$ in a Penning trap.
\end{enumerate}
Items 1--2 will be resolved within the next two to three years by DESI DR2.

%═══════════════════════════════════════════════════════════════════════════════
\section{Summary: how to evaluate the series}
%═══════════════════════════════════════════════════════════════════════════════

The density-feedback Hopfion series is a self-consistent, internally rigorous
attempt to derive the observed structure of particle physics and cosmology from
a single topological field configuration with one experimental input.
Paper~XII~\cite{Paper12} reduces the input to: $\TCMB$ alone, with $m_e$ following from
the profile normalisation proof. Papers~XIII and~XIV extend the framework to atomic quantum mechanics and close the electroweak residual to $0.04\%$.

It succeeds, within the accuracy of the computations performed, at reproducing:
the fine structure constant, the three charged lepton masses and their Koide
relation, the QCD coupling, the quark mass hierarchy, the CKM matrix, the CP
phase, the CMB spectral observables, the dark energy equation of state,
the Schr\"odinger equation, spin-$\tfrac{1}{2}$, the Pauli exclusion
principle, the electroweak VEV to $0.04\%$, the QCD scale, and the Higgs boson mass
to $0.13\sigma$ ($125.10\,\mathrm{GeV}$, residual $-0.01\%$).
It does not contradict any established SM or GR observation.

Its primary open vulnerabilities, for a rigorous evaluator, are: the $3+1$D
bosonization extension (not fully rigorous), the subleading WZW corrections
to $\aem$ (computed to one loop, not two), the experimental status of the
DESI and BAO predictions (awaiting DR2), and the free-electron mass
anisotropy prediction ($\Delta m/m_e = 1/\vp^8$, untested).

The most important single result, in the view of this guide, is Paper~X's
canonical quantisation. If the Hilbert space, Born rule, and Fibonacci spectrum
of Paper~X~\cite{Paper10} are accepted as the quantum theory of the condensate, then every
prediction of Papers~I--XI follows as consequences of a single, mathematically
well-defined quantum field theory. Paper~XII~\cite{Paper12} then closes the chain: the condensate
scale, and through it the electron mass, is fixed by $\TCMB$ alone via the
thermodynamic equilibrium of the condensate with the CMB photon bath.
Paper~XIII~\cite{Paper13} shows that standard quantum mechanics (the Schr\"odinger equation,
spin, and the Pauli principle) follows from the condensate without additional
postulates. Paper~XIV~\cite{Paper14} closes the electroweak residual to $0.04\%$ via the
exact toroidal correction and derives the Higgs boson mass
$\mH = \sqrt{k+2}\,r_{\mathrm{WZW}}\,\vEW = 125.10\,\mathrm{GeV}$
($0.13\sigma$ from ATLAS combined) from the condensate's scalar sector, using the
pentagon Parseval sum $(k+2)/2$ and the KZB torus conformal-block monodromy.
Paper~XV~\cite{Paper15} extends the framework to the $Q_H=3$ quark sector via
the McKay chain $2T\leftrightarrow(E_6)_1\supset\mathrm{SU}(3)_1^{\times 3}$,
derives both quark charges ($|Q_d|=1/3$, $|Q_u|=2/3$) from the trinification
conformal embedding with no Standard Model input, and proves $\lambda_3=\vp^6$
via the triple-fusion multiplicity argument and the $j=3/2$ simple-current
sector assignment (three independent proofs).
These are the claims the series ultimately makes, and they are the claims that
warrant the most careful scrutiny.

%─── References ───────────────────────────────────────────────────────────────
\begin{thebibliography}{99}

\bibitem{Paper1}
F.~Manfredi,
\textit{The Density-Feedback Faddeev--Niemi Hopfion:
  Exact Algebraic Predictions for Standard-Model Parameters},
Zenodo (2026).
\href{https://doi.org/10.5281/zenodo.19342027}{doi:10.5281/zenodo.19342027}

\bibitem{Paper2}
F.~Manfredi,
\textit{BPS Structure and Fixed-Point Theorem for the
  Density-Feedback Faddeev--Niemi Hopfion},
Zenodo (2026).
\href{https://doi.org/10.5281/zenodo.19363491}{doi:10.5281/zenodo.19363491}

\bibitem{Paper3}
F.~Manfredi,
\textit{Two Constants from One Knot: The Fine Structure Constant and
  the Linking Scale as Exact Predictions of the Icosahedral WZW Condensate},
Zenodo (2026).
\href{https://doi.org/10.5281/zenodo.19478629}{doi:10.5281/zenodo.19478629}

\bibitem{Paper4}
F.~Manfredi,
\textit{The Hopf Spoke, WZW Fermion Mass Renormalization,
  and Three- and Four-Term Formulas for the Fine Structure Constant},
Zenodo (2026).
\href{https://doi.org/10.5281/zenodo.19504729}{doi:10.5281/zenodo.19504729}

\bibitem{Paper5}
F.~Manfredi,
\textit{Colour Confinement, the QCD Scale, and Hadronic Structure
  from the Density-Feedback Hopfion},
Zenodo (2026).
\href{https://doi.org/10.5281/zenodo.19638857}{doi:10.5281/zenodo.19638857}

\bibitem{Paper6}
F.~Manfredi,
\textit{The Golden-Ratio Tower: Fermion Scale Hierarchy from the Hopfion RG Fixed Point},
Zenodo (2026).
\href{https://doi.org/10.5281/zenodo.19646945}{doi:10.5281/zenodo.19646945}

\bibitem{Paper7}
F.~Manfredi,
\textit{Emergent Gravity, Inflation, Dark Energy, and the Weak
  Equivalence Principle from the Density-Feedback Hopfion Condensate},
Zenodo (2026).
\href{https://doi.org/10.5281/zenodo.19646953}{doi:10.5281/zenodo.19646953}

\bibitem{Paper8}
F.~Manfredi,
\textit{Geometry-Dependent Bell Violations from the Density-Feedback Hopfion},
Zenodo (2026).
\href{https://doi.org/10.5281/zenodo.18737070}{doi:10.5281/zenodo.18737070}

\bibitem{Paper9}
F.~Manfredi,
\textit{Division Algebras, the Born Rule, and the Tsirelson Bound
  from the Density-Feedback Hopfion},
Zenodo (2026).
\href{https://doi.org/10.5281/zenodo.20075227}{doi:10.5281/zenodo.20075227}

\bibitem{Paper10}
F.~Manfredi,
\textit{Quantisation of the Hopfion Condensate: Hilbert Space,
  Liouville Measure, and the Born Rule},
Zenodo (2026).
\href{https://doi.org/10.5281/zenodo.20075279}{doi:10.5281/zenodo.20075279}

\bibitem{Paper11}
F.~Manfredi,
\textit{The Golden-Spiral as a Universal Mass Hierarchy Map:
  From Condensate Scale to Chemical Periodicity},
Zenodo (2026).
\href{https://doi.org/10.5281/zenodo.20173763}{doi:10.5281/zenodo.20173763}

\bibitem{Paper12}
F.~Manfredi,
\textit{Paper~XII: The Profile Normalisation Conjecture --- Proof via CMB Energy Identity and Two-Route Equivalence},
\href{https://doi.org/10.5281/zenodo.19646945}{doi:10.5281/zenodo.19646945}

\bibitem{Paper13}
F.~Manfredi,
\textit{Atomic Quantum Mechanics from the Hopfion Condensate},
Zenodo (2026).
\href{https://doi.org/10.5281/zenodo.20471482}{doi:10.5281/zenodo.20471482}

\bibitem{Paper14}
F.~Manfredi,
\textit{The Thick-Torus Profile Correction to the Density-Feedback Hopfion},
Zenodo (2026).
\href{https://doi.org/10.5281/zenodo.20479790}{doi:10.5281/zenodo.20479790}

\bibitem{Paper15}
F.~Manfredi,
\textit{The $Q_H=3$ Sector of the Density-Feedback Hopfion},
Zenodo (2026).
\href{https://doi.org/10.5281/zenodo.20691001}{doi:10.5281/zenodo.20691001}

\bibitem{Paper16}
F.~Manfredi,
\textit{Quark Generation Masses in the Density-Feedback Hopfion:
A Survey of Ruled-Out Mechanisms and the $Q_H=3$ Gradient-Flow Landscape},
Zenodo (2026).
\href{https://doi.org/10.5281/zenodo.21012650}{doi:10.5281/zenodo.21012650}

\bibitem{Paper17}
F.~Manfredi,
\textit{The $Q_H=1$ Sector of the Density-Feedback Hopfion:
Topology, Group Structure, Colour Exclusion},
Zenodo (2026).
\href{https://doi.org/10.5281/zenodo.21013412}{doi:10.5281/zenodo.21013412}

\end{thebibliography}

\end{document}